% Lösungen Einheit 2 von 2 – Spiegelung an der Winkelhalbierenden (zusammenbau v0.1)
\einheitenkopf{Einheit 2 von 2 · Spiegelung an der Winkelhalbierenden}

\erg{11}{a) $(5 | 2)$ – die Koordinaten tauschen \quad b) $A'(6 | 2)$, $B'(4 | -1)$ \quad c) $\frac{1}{3}$ \quad d) Die Spiegelung an y = x bildet den Graphen von f auf den von g ab; t steht senkrecht auf y = x und geht in sich über. Also berührt t den Graphen von g im Spiegelpunkt von Q, dem Punkt mit vertauschten Koordinaten. \quad e) $20$ FE: AA' und BB' stehen senkrecht auf y = x, sind also parallel; $|AA'| = 4\sqrt{2}$, $|BB'| = 6\sqrt{2}$; Höhe = Abstand der Mittelpunkte (2 | 2) und (4 | 4), also $2\sqrt{2}$; Fläche $\frac{4\sqrt{2} + 6\sqrt{2}}{2} \cdot 2\sqrt{2} = 20$ \quad f) Wahr: links steht die Fläche zwischen y = x und dem Graphen von g, rechts die zwischen dem Graphen von f und y = x; beide Flächen sind Spiegelbilder an y = x und gleich groß.}
\erg{12}{a) Fehler: Mia erkennt y = x nicht als Symmetrieachse. Die beiden Flächen sind Spiegelbilder und damit gleich groß. \quad b) Punkt und Bildpunkt liegen auf einer Senkrechten zu y = x gleich weit von ihr entfernt; die Mitte liegt auf y = x. Das leistet genau der Punkt mit vertauschten Koordinaten. \quad c) $0{,}5$ Tage je cm – der Kehrwert der Wachstumsgeschwindigkeit ist die Steigung der Umkehrfunktion Tage in Abhängigkeit von der Höhe \quad d) Auf g: (1 | 0), (2 | 1), (4 | 2); auf f: $(0 | 1)$, $(1 | 2)$, $(2 | 4)$}
